\section{Evaluation}

The evaluation asks whether one orchestration and evidence contract can
support heterogeneous machines while keeping observation selection, sink
costs, memory behavior, persistence, and package consumption inspectable.
Every numerical result in this section remains unresolved until the sealed
artifact receipts arrive. The source and package identities used for final
claims will be injected through GlassVM source \GlassVMSourceSHA, machine
source \MachinesSourceSHA, GlassVM package \GlassVMVersion, and machine
package \MachinesVersion; the exact setup coordinates are
\TBASourceCoordinates.

\subsection{Experimental setup}
\label{sec:setup}

\paragraph{Question.} Which exact source trees, package artifacts, hosts, and
run configurations produced the evidence reported here?

\paragraph{Method.} Re-run the prescribed evaluation matrix from clean
checkouts and sealed package artifacts. Record source revisions, package
versions, toolchains, host identity, observation requests, run controls, and
artifact digests in the sealed manifest. No intermediate summary supplies a
final result.

\paragraph{Metrics.} Final source and package coordinates, environment
identifiers, workload definitions, and repetition structure are
\TBASourceCoordinates. These values will be read from the sealed receipts.

\paragraph{Expected artifact source.} Sealed experiment manifest and source
attestations for the GlassVM and machine repositories.

\paragraph{Result.} \TBASourceCoordinates

\subsection{Shared-contract conformance}
\label{sec:conformance}

\paragraph{Question.} Can the same preparation, observation negotiation,
execution, and accounting path serve all evaluated machine bundles?

\paragraph{Method.} Exercise each provider through the shared public
contracts. Check successful preparation and execution, required-observation
prerequisites, optional capability downgrades, channel delivery, receipt
statuses, recorder finalization, and publication behavior.

\paragraph{Metrics.} Per-machine and per-contract-case outcomes, including
required and optional prerequisite handling, fulfilled channels, and
separate execution, evidence, sink, recorder, and publication statuses.

\paragraph{Expected artifact source.} Sealed conformance receipts tied to the
source coordinates in Section~\ref{sec:setup}.

\paragraph{Result.} The shared path is implemented across CHIP-8, PICO-8, and
TIC-80. The final conformance summary is \TBAConformanceResult.

\subsection{Selective observation cost}
\label{sec:observation-cost}

\paragraph{Question.} What resource cost is associated with selecting
different observation channels and capabilities?

\paragraph{Method.} Compare runs under the sealed observation modes while
holding each machine's workload and execution controls fixed. Account for
the selected channel set and requested normalizer dependencies in each run.

\paragraph{Metrics.} Execution time and resource use by observation mode,
with the measured differences and uncertainty left to the sealed result.

\paragraph{Expected artifact source.} Sealed selective-observation runs,
receipt data, and the figure generator identified by the artifact manifest.

\paragraph{Result.} \TBAObservationResult

\begin{figure}[t]
  \centering
  \PaperPlaceholderFigure{2.4cm}{\TBAObservationFigure}
  \caption{Selective observation cost across negotiated modes.
  \TBAObservationCaption}
  \label{fig:observation-cost}
\end{figure}

\subsection{Run-length memory behavior}
\label{sec:memory}

We test whether selected observation paths retain run-length-dependent state
inside GlassVM by executing geometrically increasing frame counts under
multiple observation modes. The comparison distinguishes process memory,
normalizer state, recorder buffering, and persisted evidence. Bounded
recorder buffering does not imply a bound on total retained run data.

\paragraph{Question.} How do in-process memory and retained evidence change
as execution length increases under the selected observation paths?

\paragraph{Method.} Execute the sealed workload matrix over increasing run
lengths with observation disabled, selected streaming reducers, and durable
recording as applicable. Measure the process and channel accounting without
conflating transient memory with on-disk evidence volume.

\paragraph{Metrics.} Peak resident memory, reported live recorder buffers,
logical evidence volume, and persisted bytes by run length and observation
mode. Numerical estimates and uncertainty are reserved for sealed results.

\paragraph{Expected artifact source.} Sealed memory-run records, recorder and
evidence receipts, and the sealed plot generator.

\paragraph{Result.} \TBAMemoryResult.\par
\noindent Evidence-volume result: \TBAStorageResult.

\begin{figure}[t]
  \centering
  \PaperPlaceholderFigure{2.6cm}{\TBAMemoryFigure}
  \caption{Peak resident memory and retained evidence versus run length.
  \TBAMemoryCaption}
  \label{fig:memory-scaling}
\end{figure}

\subsection{Python consumer overhead}
\label{sec:python-sink}

\paragraph{Question.} What is the cost of delivering owned emission batches
to a synchronous Python consumer?

\paragraph{Method.} Compare the sealed Rust-side execution path with Python
sink delivery for corresponding requests and workloads. The sink receives
typed values from the same logical emission stream; it does not write a
second durable format. Record callback work and execution outcomes separately.

\paragraph{Metrics.} Delivery time, execution time, number and logical size
of delivered batches, and any delivery failures, as reported by the sealed
measurement protocol.

\paragraph{Expected artifact source.} Sealed Python sink measurements and
their source, package, workload, and environment manifest.

\paragraph{Result.} \TBAPythonSinkResult

\begin{figure}[t]
  \centering
  \PaperPlaceholderFigure{2.4cm}{\TBAPythonFigure}
  \caption{Python sink delivery cost under the sealed workload.
  \TBAPythonCaption}
  \label{fig:python-sink-cost}
\end{figure}

\subsection{Rust/Python semantic equivalence}
\label{sec:equivalence}

The equivalence check uses one \texttt{PreparedRun}, one
\texttt{PreparedObservation}, one execution, and one emission stream routed
to a Python sink and the Rust recorder. The comparison is performed over
logical evidence per compared channel, preserving order within each channel.
It does not assert a global ordering between independent channels.

\paragraph{Question.} Do both consumers receive equivalent identities and
logical records when they consume the same prepared execution?

\paragraph{Method.} Compare \texttt{PreparedRunIdentity} and
\texttt{PreparedObservationIdentity}, then compare logical evidence for each
included channel, capability outputs, and capability receipts. Keep delivery
and persistence accounting distinct from the evidence comparison.

\paragraph{Metrics.} Identity agreement, per-channel record agreement,
capability output and receipt agreement, and the sealed summary of any
permitted representation normalization.

\paragraph{Expected artifact source.} Sealed semantic-equivalence report,
including exact source coordinates and comparison artifact digests.

\paragraph{Result.} Structural semantic equivalence has been checked across
the three machine implementations. The final sealed summary is
\TBAEquivalenceResult.\par
\noindent Exact artifact coordinates: \TBASourceCoordinates.

\subsection{Package-index-only consumption}
\label{sec:package-consumption}

\paragraph{Question.} Can an external consumer install and use the facade
and machine provider from package indexes without a local GlassVM source
checkout on its import or execution path?

\paragraph{Method.} In a clean consumer environment, install the facade and
provider from package indexes, discover the provider through its entry point,
prepare an artifact, execute it, receive an evidence receipt, finalize the
reference recorder, and verify a published run. A separate Rust consumer
resolves the published GlassVM crates from crates.io without a local source
path.

\paragraph{Metrics.} Successful package resolution, provider discovery,
preparation and execution outcomes, evidence receipt, recorder finalization,
publication reference, and the exact sealed package/source coordinates.

\paragraph{Expected artifact source.} Sealed Python and Rust external-consumer
proof manifests and package-index receipts.

\paragraph{Result.} Package-index-only consumption has been demonstrated;
final package versions and exact proof coordinates remain
\TBAPackageProofCoordinates.
